% ==============================================================================
% ANEXOS
% ==============================================================================
% Datos tomados de los repositorios del proyecto:
%   - docs/model-provenance.md (logSguarDian): hashes de los pares de modelos.
%   - SHA-256 de los archivos .tgz de npm: calculados con npm pack y shasum el
%     2026-10-12.
%   - docs/cybersecurity-objectives-compliance.md: historial del bloqueo del
%     conjunto de prueba (git log --all -- training/splits/test.lock.sha256).
%   - docs/dataset-audit.md: auditoria de fuentes.
%   - docs/evidence/README.md (logSguarDian-vulnerable-project): corrida de
%     latencia en Linux.
\ifdefined\modoensamble
\else
	\documentclass[11pt, letterpaper]{report}
\usepackage[spanish, es-nodecimaldot, es-noquoting]{babel}
\selectlanguage{spanish}
\usepackage[utf8]{inputenc}
\usepackage[T1]{fontenc}

% ===============================================================================
% El estudiante debe llenar sus datos en esta sección para que la plantilla los 
% auto-importe y genere automáticamente las páginas de portada y de firmas 
% autorizadas.
% ===============================================================================
% Datos del estudiante:
% -------------------------------------------------------------------------------
% Nombre completo
\def \nombreestudiante {Sebastián Huertas Gómez}
% Carné
\def \uvgcarne {22295}
% Facultad
\def \uvgfacultad {Ingeniería}
% Carrera
\def \uvgcarrera {Ingeniería en Ciencia de la Computación y Tecnologías de la Información}

% Datos del trabajo:
% -------------------------------------------------------------------------------
% Título completo (tomado literalmente de protocolo_sebastian_v1 (3).pdf, sección 1)
\def \titulotesis {logSguarDian: Desarrollo de una librería para la detección de amenazas de seguridad en Endpoints de aplicaciones}
% Mes y año de entrega (confirmado por el estudiante: noviembre 2026)
\def \mesentrega {noviembre}
\def \anoentrega {2026}
% Asesor (según protocolo_sebastian_v1 (3).pdf, portada)
\def \nombreasesor {Ing. Erick Francisco Marroquín Rodríguez}
% El catedrático del curso (Ing. Bacilio Alexander Bolaños Lima) figura como
% primer examinador en la hoja de firmas, definido más abajo.

% Datos del tribunal examinador:
% -------------------------------------------------------------------------------
% Nombre del primer examinador (segunda línea de la hoja de firmas)
\def \nombreprimerex {Ing. Ma. Bacilio Alexander Bolaños Lima}

% Capítulos pre-definidos
% -------------------------------------------------------------------------------
% Comentar las líneas de las secciones que desean omitirse, por defecto se
% se incluyen todas.
% Se descomentan conforme se redacta cada capítulo (todavía sin archivo de
% contenido: prefacio, alcance, anexos, glosario).
\def \CAPprefacio {Prefacio}
\def \CAPantecedentes {Antecedentes}
\def \CAPalcance {Alcance}
\def \CAPmetodologia {Metodología}
\def \CAPresultados {Resultados}
\def \CAPdiscusion {Análisis de resultados}
\def \CAPanexos {Anexos}
\def \CAPglosario {Glosario}

% Formato y estilo de la plantilla
% -------------------------------------------------------------------------------
% Modo impresión: Puede des-comentar la siguiente línea para generar un documento pdf sin la portada, para cuando se desee imprimir el documento para encuadernación
% NOTA: la plantilla original traía esta línea activa por defecto, lo cual
% suprime la portada verde (main.tex hace \let\CAPportada\undefined cuando
% \printver está definido). Se comenta para que la versión digital, con
% portada, sea la que se genera por defecto.
%\def \printver {Versión del documento para impresión}

% Portada: Puede cambiarse la imagen en la portada al cambiar el nombre del 
% archivo siguiente. NOTA: debe tener la suficiente resolución para cubrir el área
% designada
\def \imagenportada {plantilla/portadacit.jpg}

% Referencias: Puede des-comentar la siguiente línea para utilizar el formato de referencias APA
%\def \usarAPA {Usar formato APA}

% Párrafo: Puede comentar la siguiente línea si desea emplear un formato de 
% párrafo distinto al establecido por defecto
\def \parpordefecto {Formato de párrafo por defecto}

% Capítulos y secciones: Puede des-comentar la siguiente línea para establecer el
% formato de los capítulos y secciones bajo el estándar original de UVG para
% trabajos de graduación. Este incluye: capítulos con numeración romana, secciones
% con letras mayúsculas, sub-secciones con números y sub-sub-secciones con letras
% minúsculas
%\def \capsecuvg {Formato UVG para capítulos y secciones}
% ================================================================================
% En este archivo se colocan opciones adicionales para modificar el formato de la
% plantilla, para emplearse en otros tipos de documentos que no sean trabajos de
% graduación. Si usted está trabajando su tesis, NO modifique este archivo
% ================================================================================
% Capítulos pre-definidos
% --------------------------------------------------------------------------------
% Comentar las líneas de las secciones que desean omitirse, por defecto se 
% se incluyen todas.
\def \CAPportada {Portada}
\def \CAPcaratula {Caratula}
\def \CAPfirmas {Hoja de firmas}
\def \CAPindice {Índice general}
\def \CAPfiguras {Listado de figuras}
\def \CAPcuadros {Listado de cuadros}
\def \CAPresumen {Resumen}
\def \CAPabstract {Resumen}
\def \CAPintroduccion {Introducción}
\def \CAPobjetivos {Objetivos}
\def \CAPjustificacion {Justificación}
\def \CAPmarcoteorico {Marco teórico}
\def \CAPconclusiones {Conclusiones}
\def \CAPrecomendaciones {Recomendaciones}
\def \CAPbibliografia {Bibliografía}

\usepackage{xcolor}
\usepackage{amsfonts}
\usepackage{amsmath}
\usepackage{amssymb}
\usepackage{amsthm}
\usepackage{mathtools}
\usepackage{graphicx}
\usepackage{xfrac}
\usepackage{float}
\usepackage[hypertexnames=false]{hyperref}
\usepackage[font=small]{caption}
\usepackage{subcaption}
% Cuadros: array (columnas L{} alineadas a la izquierda) y longtable (cuadros
% que ocupan más de una página). Se cargan aquí, no en
% 2-paquetes_y_comandos_usuario.tex, para que los capítulos de Metodología y
% Resultados no dependan de ese archivo.
\usepackage{array}
\usepackage{longtable}
\newcolumntype{L}[1]{>{\raggedright\arraybackslash}p{#1}}
\renewcommand{\arraystretch}{1.3}
\setlength{\LTpre}{0.3cm}
\setlength{\LTpost}{0.1cm}
% Margen de estiramiento para que los párrafos con palabras largas no
% generen líneas que sobresalgan del margen.
\setlength{\emergencystretch}{3em}
\usepackage{xpatch}
\usepackage{emptypage}
% htt: permite dividir con guion el texto en \texttt (nombres de archivo,
% clases y campos), que de otro modo sobresale del margen derecho.
\usepackage[htt]{hyphenat}
\usepackage{fancyhdr}
\usepackage[backend=biber, style=ieee]{biblatex}
\ifdefined\usarAPA
    \usepackage[backend=biber, style=apa]{biblatex}
\fi
\addbibresource{m-bibliografia.bib}

\usepackage[percent]{overpic}
\usepackage{chngcntr}

\ifdefined\CAPglosario
	\usepackage[numberedsection]{glossaries}
	\makeglossaries
    \newglossaryentry{rasp}
{
    name={RASP},
    description={Autoprotección de aplicaciones en tiempo de ejecución (\textit{Runtime Application Self-Protection}): mecanismo de seguridad que se ejecuta dentro de la propia aplicación y analiza sus peticiones para detectar y bloquear ataques.}
}

\newglossaryentry{waf}
{
    name={WAF},
    description={Cortafuegos de aplicaciones web (\textit{Web Application Firewall}): dispositivo o servicio perimetral que filtra el tráfico HTTP antes de que llegue a la aplicación.}
}

\newglossaryentry{modsecurity}
{
    name={ModSecurity},
    description={Motor de WAF de código abierto que se usa con reglas externas; en este trabajo se configuró con OWASP CRS.}
}

\newglossaryentry{crs}
{
    name={OWASP CRS},
    description={Conjunto de reglas genéricas de OWASP (\textit{Core Rule Set}) para ModSecurity, con niveles de paranoia que controlan qué tan estrictas son.}
}

\newglossaryentry{owasp}
{
    name={OWASP Top 10},
    description={Lista de las diez categorías de riesgo más críticas para aplicaciones web, publicada por OWASP; este trabajo usa la edición 2021.}
}

\newglossaryentry{mitre}
{
    name={MITRE ATT\&CK},
    sort={MITRE ATT&CK},
    description={Base de conocimiento de tácticas y técnicas de atacantes, usada aquí para asociar cada vector de ataque a una técnica.}
}

\newglossaryentry{cwe}
{
    name={CWE},
    description={Enumeración común de debilidades de software (\textit{Common Weakness Enumeration}): catálogo de tipos de defectos que originan vulnerabilidades.}
}

\newglossaryentry{sqli}
{
    name={SQL Injection},
    description={Ataque que inserta sentencias SQL en una entrada de la aplicación para alterar las consultas que ejecuta la base de datos.}
}

\newglossaryentry{xss}
{
    name={XSS},
    description={\textit{Cross-Site Scripting}: ataque que inyecta código ejecutable en una página que otros usuarios verán.}
}

\newglossaryentry{pathtraversal}
{
    name={Path Traversal},
    description={Ataque que manipula rutas de archivo (por ejemplo con \texttt{../}) para leer archivos fuera del directorio previsto; incluye la inclusión local de archivos (LFI).}
}

\newglossaryentry{cmdi}
{
    name={Command Injection},
    description={Ataque que logra ejecutar comandos del sistema operativo a través de una entrada de la aplicación.}
}

\newglossaryentry{middleware}
{
    name={\textit{middleware}},
    sort={middleware},
    description={Función que se ejecuta entre la recepción de una petición y la lógica de negocio de la aplicación; en Express se registra con \texttt{app.use}.}
}

\newglossaryentry{rf}
{
    name={Random Forest},
    description={Algoritmo supervisado que combina muchos árboles de decisión; en la librería es el único modelo con autoridad para bloquear una petición.}
}

\newglossaryentry{iforest}
{
    name={Isolation Forest},
    description={Algoritmo no supervisado de detección de anomalías que aísla las observaciones atípicas con particiones aleatorias; en la librería solo registra anomalías y nunca bloquea.}
}

\newglossaryentry{onnx}
{
    name={ONNX},
    description={Formato abierto para serializar modelos de aprendizaje automático, que permite ejecutarlos fuera de Python.}
}

\newglossaryentry{workerthreads}
{
    name={\texttt{worker\_threads}},
    sort={worker_threads},
    description={Módulo de Node.js que ejecuta código en hilos independientes del hilo principal; la inferencia de los modelos se ejecuta ahí.}
}

\newglossaryentry{eventloop}
{
    name={\textit{Event Loop}},
    sort={Event Loop},
    description={Bucle de eventos de Node.js que atiende las peticiones en un único hilo; bloquearlo degrada a toda la aplicación.}
}

\newglossaryentry{failopen}
{
    name={\textit{fail-open}},
    sort={fail-open},
    description={Política por la cual, ante un error o un tiempo de espera del módulo de detección, la petición se deja pasar en lugar de bloquearse.}
}

\newglossaryentry{monitor}
{
    name={Modo \texttt{monitor}},
    sort={Modo monitor},
    description={Modo de operación en el que la librería registra las peticiones que bloquearía, pero no las bloquea.}
}

\newglossaryentry{block}
{
    name={Modo \texttt{block}},
    sort={Modo block},
    description={Modo de operación en el que la librería responde HTTP 403 a las peticiones que clasifica como ataque.}
}

\newglossaryentry{pesp}
{
    name={p95},
    description={Percentil 95 del tiempo de respuesta: el valor por debajo del cual queda el 95~\% de las peticiones.}
}

\newglossaryentry{deltap95}
{
    name={$\Delta p95$},
    sort={Delta p95},
    description={Diferencia entre el p95 con el middleware activo y el p95 de la línea base sin él; en este trabajo se mide en milisegundos.}
}

\newglossaryentry{funo}
{
    name={F1},
    description={Media armónica de la precisión y el recall; es la métrica principal de clasificación del trabajo.}
}

\newglossaryentry{auc}
{
    name={AUC-ROC},
    description={Área bajo la curva ROC; mide qué tan bien separa un modelo dos clases sin depender de un umbral.}
}

\newglossaryentry{wilson}
{
    name={Intervalo de Wilson},
    description={Intervalo de confianza para una proporción, adecuado para muestras pequeñas o proporciones cercanas a 0 o 1.}
}

\newglossaryentry{casidup}
{
    name={Casi-duplicado},
    description={Registro muy parecido a otro pero no idéntico; si queda repartido entre entrenamiento y prueba, puede inflar las métricas.}
}

\newglossaryentry{minhash}
{
    name={MinHash},
    description={Técnica para estimar rápidamente la similitud entre conjuntos, usada para detectar casi-duplicados.}
}

\newglossaryentry{fuga}
{
    name={Fuga de datos},
    description={Presencia en el conjunto de evaluación de información disponible durante el entrenamiento, lo que sobreestima el desempeño.}
}

\newglossaryentry{corpuslimpio}
{
    name={Corpus limpio},
    description={Conjunto de payloads de evaluación filtrado para excluir los que coinciden con el entrenamiento; en este trabajo es el de la Ronda~5.}
}

\newglossaryentry{useragent}
{
    name={\texttt{User-Agent}},
    sort={User-Agent},
    description={Cabecera HTTP con la que el cliente se identifica (navegador, cURL, entre otros); su representación en el corpus sesgó al modelo.}
}

\newglossaryentry{falsopositivo}
{
    name={Falso positivo},
    description={Petición legítima que el sistema clasifica o bloquea como ataque.}
}

\newglossaryentry{parmodelos}
{
    name={Par de modelos},
    description={Combinación de un archivo \texttt{rf.onnx} y uno \texttt{if.onnx} identificada por su suma SHA-256 (pares A, B y C).}
}

\newglossaryentry{tarball}
{
    name={\textit{tarball}},
    sort={tarball},
    description={Archivo comprimido \texttt{.tgz} de un paquete de npm, publicado o construido localmente.}
}

\newglossaryentry{kanban}
{
    name={Kanban},
    description={Metodología de gestión basada en un tablero de columnas que representan el estado de las tareas.}
}

\newglossaryentry{artillery}
{
    name={Artillery},
    description={Herramienta de pruebas de carga usada para medir la latencia con y sin el middleware.}
}

\newglossaryentry{sqlmap}
{
    name={sqlmap},
    description={Herramienta de código abierto que automatiza la detección y la explotación de inyecciones SQL; generó los payloads de SQL Injection del corpus limpio.}
}

\newglossaryentry{zap}
{
    name={OWASP ZAP},
    description={Herramienta de OWASP para pruebas de seguridad de aplicaciones web; se usó para confirmar el payload de XSS.}
}

\newglossaryentry{mcp}
{
    name={MCP},
    description={Protocolo de Contexto de Modelo (\textit{Model Context Protocol}): estándar para que un asistente de inteligencia artificial use herramientas externas.}
}

\fi

\usepackage[top=1in, left=1.5in, right=1in, bottom=1in]{geometry}
\usepackage[Sonny]{fncychap}

\graphicspath{ {figuras/} }
\definecolor{uvg-green}{RGB}{17,71,52}
\newcommand{\defaultparformat}[1]{
	{\setlength{\parskip}{2ex}
     \input{#1}}
}
\ifdefined\capsecuvg
	\renewcommand\thechapter{\Roman{chapter}}
    \renewcommand\thesection{\Alph{section}}
	\renewcommand\thesubsection{\arabic{subsection}}
    \renewcommand\thesubsubsection{\alph{subsubection}}
\fi
\counterwithout{figure}{chapter}
\counterwithout{table}{chapter}
\counterwithout{equation}{chapter}

\providecommand{\blankpage}{
\newpage
\thispagestyle{empty}
\mbox{}
\newpage
}

% ==============================================================================
% Paquetes y comandos adicionales del usuario
% ==============================================================================
% TikZ: requerido por las figuras del capítulo de Marco Teórico (pipeline de
% middleware, comparación RASP vs. control perimetral).
\usepackage{tikz}
\usetikzlibrary{arrows.meta, positioning, shapes.geometric, fit, backgrounds}


\pagestyle{headings}
\begin{document}

\ifdefined\parpordefecto
	\setlength{\parskip}{2ex}
\fi

\fi

\chapter{Anexos}
\label{cap:anexos}

\section{Procedencia de los modelos y de los paquetes}
\label{anx:procedencia}

El Cuadro~\ref{tab:anx-hashes} lista los archivos \texttt{rf.onnx} e \texttt{if.onnx} de cada par
de modelos descrito en el Cuadro~\ref{tab:pares-modelos}, con su suma SHA-256 completa. Las sumas
se calcularon sobre el contenido del objeto de git de cada archivo y, para los paquetes, sobre el
contenido extraído del archivo publicado en npm.

\begin{table}[H]
\centering
\small
\begin{tabular}{|L{3.6cm}|L{9.9cm}|}
\hline
\textbf{Archivo} & \textbf{SHA-256} \\ \hline
Par A, \texttt{rf.onnx} (\texttt{rf\_v10}) & {\scriptsize\texttt{180837255dffd596a38bf7c29032ccc3\allowbreak{}13005cd70eca5c6238e23e35bdc0d4ec}} \\ \hline
Par A, \texttt{if.onnx} (\texttt{if\_v9}) & {\scriptsize\texttt{5de374a7f1d19176d7c66d3354e972d3\allowbreak{}7b28a7306a87efa0ce8639278f489150}} \\ \hline
Par B, \texttt{rf.onnx} (\texttt{rf\_v11}, primer build) & {\scriptsize\texttt{a7c015f5a5fa3fabfd91b79ec9338dc7\allowbreak{}da6a68c392c19ded4b04cc6f0fb903b2}} \\ \hline
Par B, \texttt{if.onnx} (\texttt{if\_v10}, primer build) & {\scriptsize\texttt{3cf8e45c285de19bc86e6dc284e502cd\allowbreak{}2c24becb888747b9645bed4cfb224417}} \\ \hline
Par C, \texttt{rf.onnx} (\texttt{rf\_v11}, reentrenamiento) & {\scriptsize\texttt{25b407e649283606a231665a5a8c8ea1\allowbreak{}c2c7ad4fd0374a8945b3f119651b0870}} \\ \hline
Par C, \texttt{if.onnx} (\texttt{if\_v10}, reentrenamiento) & {\scriptsize\texttt{5be28ee8437643b2f13708eee1ac3175\allowbreak{}348a350ea6fb8f8daeb09bd5717f797e}} \\ \hline
\end{tabular}
\caption[Sumas SHA-256 de los pares de modelos]{Sumas SHA-256 de los archivos ONNX de cada par de modelos}
\label{tab:anx-hashes}
\end{table}
{\footnotesize Nota. Elaboración propia con datos de \texttt{docs/model-provenance.md} (repositorio
\texttt{logSguarDian}). El par~A es el contenido del paquete \texttt{logsguardian@0.1.0}; el par~C,
el del paquete \texttt{logsguardian@0.1.1}.}

El Cuadro~\ref{tab:anx-paquetes} reúne los dos paquetes publicados en el registro de npm y el umbral
del Isolation Forest que cada uno aplica.

\begin{table}[H]
\centering
\small
\begin{tabular}{|L{2.3cm}|L{2.5cm}|L{6.2cm}|L{2.5cm}|}
\hline
\textbf{Paquete} & \textbf{Publicación} & \textbf{SHA-256 del archivo \texttt{.tgz}} & \textbf{Umbral del IF} \\ \hline
\texttt{0.1.0} & 30 de agosto de 2026 & {\scriptsize\texttt{c5c1a1da57cfb873350ac0c8ab3e7bad\allowbreak{}395e8cbf507b18cc3272df447880dcb0}} & 0.002486 (constante) \\ \hline
\texttt{0.1.1} & 11 de octubre de 2026 (UTC) & {\scriptsize\texttt{abf8d5c47036d9c3a289bc36158f1caa\allowbreak{}a91b9504689aa97193e51c6808258d15}} & 0.004206 (leído de \texttt{parity\_report.json}) \\ \hline
\end{tabular}
\caption[Paquetes publicados en npm]{Paquetes \texttt{logsguardian} publicados en el registro de npm}
\label{tab:anx-paquetes}
\end{table}
{\footnotesize Nota. Elaboración propia. Las sumas de los archivos se calcularon descargando cada
versión del registro (\texttt{npm pack}); los umbrales provienen de \texttt{docs/decision-policy.md}
del repositorio \texttt{logSguarDian}. Los otros dos archivos que se mencionan en el texto, el
tarball de la Ronda~5 original (\texttt{e09120006dd1...}) y el tarball de validación
(\texttt{cf1238deb333...}), no se publicaron en el registro.}

\section{Historial del bloqueo del conjunto de prueba}
\label{anx:lock}

El bloqueo del conjunto de prueba es el SHA-256 de los identificadores por fila del conjunto,
ordenados y unidos por saltos de línea. El valor cambió nueve veces; el Cuadro~\ref{tab:anx-lock}
las lista. No hay un único bloqueo por generación de modelo, y el cambio del 23 de septiembre de
2026 entró en tres commits.

\begin{table}[H]
\centering
\small
\begin{tabular}{|L{2.2cm}|L{2.4cm}|L{4.8cm}|L{2.2cm}|L{2.2cm}|}
\hline
\textbf{Fecha} & \textbf{Commit} & \textbf{Motivo} & \textbf{Antes} & \textbf{Después} \\ \hline
11 jun 2026 & \texttt{cc00c3c} & Bloqueo inicial & --- & \texttt{3beeaba4} \\ \hline
23 jul 2026 & \texttt{df62619} & Nuevo conjunto, antes de la generación v3 & \texttt{3beeaba4} & \texttt{d1a2a284} \\ \hline
24 jul 2026 & \texttt{6867f28} & Antes de la generación v4 & \texttt{d1a2a284} & \texttt{7df5adb8} \\ \hline
24 jul 2026 & \texttt{59990ca} & Antes de la generación v6 & \texttt{7df5adb8} & \texttt{ee4b85db} \\ \hline
25 jul 2026 & \texttt{5dc7c06} & Antes de la generación v8 & \texttt{ee4b85db} & \texttt{251716bb} \\ \hline
26 jul 2026 & \texttt{fa46bdd} & Antes de la generación v9 & \texttt{251716bb} & \texttt{4a5e0b97} \\ \hline
2 ago 2026 & \texttt{b43de9c} & Antes de la generación v10 & \texttt{4a5e0b97} & \texttt{6c5f253d} \\ \hline
30 ago 2026 & \texttt{8cbe6c9} & Antes de la generación v11 & \texttt{6c5f253d} & \texttt{a3140ab3} \\ \hline
23 sep 2026 & \texttt{f15e7cc}, \texttt{72aabd6}, \texttt{6d1fa68} & Rebloqueo para que coincida con la generación v11 restaurada & \texttt{a3140ab3} & \texttt{7486c258} \\ \hline
\end{tabular}
\caption[Historial del bloqueo del conjunto de prueba]{Historial de cambios del bloqueo del conjunto de prueba (hashes truncados a 8 caracteres)}
\label{tab:anx-lock}
\end{table}
{\footnotesize Nota. Elaboración propia con datos de \texttt{docs/cybersecurity-objectives-compliance.md}
(repositorio \texttt{logSguarDian}), generados con \texttt{git log --all -{}- training/splits/test.lock.sha256}.}

\section{Auditoría de las fuentes de datos}
\label{anx:auditoria}

El Cuadro~\ref{tab:anx-auditoria} resume la auditoría de licencias de las nueve fuentes públicas
candidatas, verificada contra la API o la página de cada origen.

\begin{table}[H]
\centering
\small
\begin{tabular}{|L{3.8cm}|L{3.0cm}|L{3.3cm}|L{3.6cm}|}
\hline
\textbf{Fuente} & \textbf{Licencia} & \textbf{DOI} & \textbf{Decisión} \\ \hline
Colección CAPEC (SR-BH 2020) \cite{sureda2022srbh,sureda2022newmultilabel} & CC0-1.0 & 10.7910/DVN/OGOIXX & Incluida \\ \hline
modsec-learn \cite{pralab2024modseclearn,scano2024modseclearn} & MIT & No verificado & Incluida \\ \hline
OWASP ModSecurity, registros de un servidor protegido (Zenodo) \cite{lucz2025} & CC BY 4.0 & 10.5281/zenodo.17178461 & Incluida \\ \hline
Payloads.csv (Kaggle) \cite{shahane2021xss} & CC BY 4.0 & No asignado por Kaggle & Incluida \\ \hline
payload\_full.csv (Kaggle) \cite{thakur2025payloads} & MIT & No asignado por Kaggle & Incluida \\ \hline
XSS\_dataset.csv (Kaggle) \cite{shah2020xss} & No especificada por el origen & No asignado por Kaggle & Incluida; no se redistribuye \\ \hline
command injection.csv (Kaggle) \cite{pawase2024cmdi} & No especificada por el origen & No asignado por Kaggle & Incluida; no se redistribuye \\ \hline
Lista de payloads de Path Traversal de omurugur \cite{ugur2019traversal} & MIT & No asignado & Aprobada; sin analizador, no aporta filas \\ \hline
AIT-LDSv2 (russellmitchell) \cite{landauer2022ait,landauer2023maintainable} & CC-BY-NC-SA-4.0 & 10.5281/zenodo.5789064 & Excluida: sus etiquetas describen el rol del atacante, no el tipo de payload \\ \hline
\end{tabular}
\caption[Auditoría de las fuentes de datos]{Licencia, DOI y decisión de inclusión de cada fuente auditada}
\label{tab:anx-auditoria}
\end{table}
{\footnotesize Nota. Elaboración propia con datos de \texttt{docs/dataset-audit.md} (repositorio
\texttt{logSguarDian}). El número de filas que cada fuente aporta al corpus final está en el
Cuadro~\ref{tab:met-origenes}. Las listas SecLists y PayloadsAllTheThings no constan en esta
auditoría; su licencia (MIT) se consultó en sus repositorios el 10 de octubre de 2026.}

\section{Evidencia de la corrida de latencia en Linux}
\label{anx:latencia}

El Cuadro~\ref{tab:anx-latencia} identifica la corrida de integración continua de la que provienen
las cifras de latencia, memoria y CPU del Capítulo~\ref{cap:resultados}.

\begin{table}[H]
\centering
\small
\begin{tabular}{|L{4.2cm}|L{9.3cm}|}
\hline
\textbf{Campo} & \textbf{Valor} \\ \hline
Flujo de trabajo & \texttt{.github/workflows/latency-benchmark.yml} \\ \hline
Identificador de la ejecución & 36276690236 \\ \hline
Identificador del artefacto & 10917957439 \\ \hline
Commit & \texttt{f15a9e8} \\ \hline
Fecha & 26 de septiembre de 2026 \\ \hline
Ejecutor & Ubuntu (Azure), x86\_64, 2 vCPU, 7.8~GiB de memoria \\ \hline
Docker Engine / Compose & 28.0.4 / v2.38.2 \\ \hline
Modelos & Par~B (\texttt{rf.onnx} \texttt{a7c015f5...}, \texttt{if.onnx} \texttt{3cf8e45c...}) \\ \hline
Contenido del archivo & 30 reportes de Artillery (3 escenarios, 2 condiciones, 5 repeticiones), muestras de \texttt{docker stats} cada $\approx$5~s y análisis de memoria \\ \hline
Copia en el repositorio & \texttt{docs/evidence/latency-linux-2026-09-26.tgz} (SHA-256 {\scriptsize\texttt{b96d901408547331fb4859a5dac4810a\allowbreak{}bdb1caaeb2f15f3378877c2290187cec}}) \\ \hline
\end{tabular}
\caption[Corrida de latencia en Linux]{Identificación de la corrida de latencia, memoria y CPU en Linux nativo}
\label{tab:anx-latencia}
\end{table}
{\footnotesize Nota. Elaboración propia con datos de \texttt{docs/evidence/README.md} (repositorio
\texttt{logSguarDian-vulnerable-project}). El artefacto original de GitHub Actions expira el 26 de
octubre de 2026; la copia del repositorio es la que se conserva.}

\section{Cobertura de XSS por estrato}
\label{anx:xss-estratos}

El Cuadro~\ref{tab:anx-xss-estratos} desglosa la cobertura del corpus propio de XSS por cada
dimensión de su generador, y el Cuadro~\ref{tab:anx-xss-bio} repite el desglose para los 64
payloads dirigidos al campo \texttt{bio}, el único con vulnerabilidad confirmada. Las
Configuraciones 2 y 3b bloquearon todos los payloads de cada estrato. La Configuración 3a falló en
11 payloads, todos con doble codificación URL (\texttt{url\_double}, 14 de 25 bloqueados) y
enviados como JSON o \texttt{multipart/form-data}. Los intervalos son de Wilson al 95~\%; en un
estrato con pocos payloads, el límite inferior es amplio aun con 100~\% de bloqueo.

\begin{longtable}{|L{4.2cm}|L{0.9cm}|L{3.4cm}|L{4.0cm}|}
\hline
\textbf{Estrato} & \textbf{n} & \textbf{Configuraciones 2 y 3b: bloqueados e IC 95~\%} & \textbf{Configuración 3a: bloqueados e IC 95~\%} \\ \hline
\endfirsthead
\multicolumn{4}{l}{\footnotesize\textit{(continuación)}} \\ \hline
\textbf{Estrato} & \textbf{n} & \textbf{Configuraciones 2 y 3b: bloqueados e IC 95~\%} & \textbf{Configuración 3a: bloqueados e IC 95~\%} \\ \hline
\endhead
\multicolumn{4}{r}{\footnotesize\textit{Continúa en la página siguiente}} \\
\endfoot
\endlastfoot
\multicolumn{4}{|l|}{\textit{Contexto de inyección}} \\ \hline
\texttt{attr\_double} & 41 & 41/41 (100.0~\%) [91.4--100.0] & 39/41 (95.1~\%) [83.9--98.7] \\ \hline
\texttt{attr\_single} & 44 & 44/44 (100.0~\%) [92.0--100.0] & 41/44 (93.2~\%) [81.8--97.7] \\ \hline
\texttt{attr\_unquoted} & 23 & 23/23 (100.0~\%) [85.7--100.0] & 22/23 (95.7~\%) [79.0--99.2] \\ \hline
\texttt{html\_body} & 47 & 47/47 (100.0~\%) [92.4--100.0] & 44/47 (93.6~\%) [82.8--97.8] \\ \hline
\texttt{html\_comment} & 16 & 16/16 (100.0~\%) [80.6--100.0] & 15/16 (93.8~\%) [71.7--98.9] \\ \hline
\texttt{js\_string} & 7 & 7/7 (100.0~\%) [64.6--100.0] & 7/7 (100.0~\%) [64.6--100.0] \\ \hline
\texttt{svg} & 13 & 13/13 (100.0~\%) [77.2--100.0] & 12/13 (92.3~\%) [66.7--98.6] \\ \hline
\texttt{url\_href\_src} & 9 & 9/9 (100.0~\%) [70.1--100.0] & 9/9 (100.0~\%) [70.1--100.0] \\ \hline
\multicolumn{4}{|l|}{\textit{Vector de ejecución}} \\ \hline
\texttt{event\_onerror} & 37 & 37/37 (100.0~\%) [90.6--100.0] & 34/37 (91.9~\%) [78.7--97.2] \\ \hline
\texttt{event\_onfocus\_autofocus} & 15 & 15/15 (100.0~\%) [79.6--100.0] & 14/15 (93.3~\%) [70.2--98.8] \\ \hline
\texttt{event\_onload\_body} & 15 & 15/15 (100.0~\%) [79.6--100.0] & 14/15 (93.3~\%) [70.2--98.8] \\ \hline
\texttt{event\_onmouseover} & 36 & 36/36 (100.0~\%) [90.4--100.0] & 33/36 (91.7~\%) [78.2--97.1] \\ \hline
\texttt{event\_ontoggle} & 25 & 25/25 (100.0~\%) [86.7--100.0] & 23/25 (92.0~\%) [75.0--97.8] \\ \hline
\texttt{iframe\_srcdoc} & 6 & 6/6 (100.0~\%) [61.0--100.0] & 6/6 (100.0~\%) [61.0--100.0] \\ \hline
\texttt{javascript\_uri} & 9 & 9/9 (100.0~\%) [70.1--100.0] & 9/9 (100.0~\%) [70.1--100.0] \\ \hline
\texttt{script} & 21 & 21/21 (100.0~\%) [84.5--100.0] & 21/21 (100.0~\%) [84.5--100.0] \\ \hline
\texttt{svg\_onload} & 36 & 36/36 (100.0~\%) [90.4--100.0] & 35/36 (97.2~\%) [85.8--99.5] \\ \hline
\multicolumn{4}{|l|}{\textit{Función ejecutada}} \\ \hline
\texttt{alert} & 28 & 28/28 (100.0~\%) [87.9--100.0] & 23/28 (82.1~\%) [64.4--92.1] \\ \hline
\texttt{confirm} & 31 & 31/31 (100.0~\%) [89.0--100.0] & 29/31 (93.5~\%) [79.3--98.2] \\ \hline
\texttt{document\_cookie} & 18 & 18/18 (100.0~\%) [82.4--100.0] & 18/18 (100.0~\%) [82.4--100.0] \\ \hline
\texttt{eval\_atob} & 32 & 32/32 (100.0~\%) [89.3--100.0] & 32/32 (100.0~\%) [89.3--100.0] \\ \hline
\texttt{fetch\_externo} & 40 & 40/40 (100.0~\%) [91.2--100.0] & 38/40 (95.0~\%) [83.5--98.6] \\ \hline
\texttt{fromcharcode} & 29 & 29/29 (100.0~\%) [88.3--100.0] & 29/29 (100.0~\%) [88.3--100.0] \\ \hline
\texttt{prompt} & 22 & 22/22 (100.0~\%) [85.1--100.0] & 20/22 (90.9~\%) [72.2--97.5] \\ \hline
\multicolumn{4}{|l|}{\textit{Codificación o evasión}} \\ \hline
\texttt{html\_entities\_dec} & 23 & 23/23 (100.0~\%) [85.7--100.0] & 23/23 (100.0~\%) [85.7--100.0] \\ \hline
\texttt{html\_entities\_hex} & 21 & 21/21 (100.0~\%) [84.5--100.0] & 21/21 (100.0~\%) [84.5--100.0] \\ \hline
\texttt{inline\_comment} & 23 & 23/23 (100.0~\%) [85.7--100.0] & 23/23 (100.0~\%) [85.7--100.0] \\ \hline
\texttt{mixed\_case} & 15 & 15/15 (100.0~\%) [79.6--100.0] & 15/15 (100.0~\%) [79.6--100.0] \\ \hline
\texttt{raw} & 14 & 14/14 (100.0~\%) [78.5--100.0] & 14/14 (100.0~\%) [78.5--100.0] \\ \hline
\texttt{slash\_for\_space} & 16 & 16/16 (100.0~\%) [80.6--100.0] & 16/16 (100.0~\%) [80.6--100.0] \\ \hline
\texttt{unicode\_escape} & 18 & 18/18 (100.0~\%) [82.4--100.0] & 18/18 (100.0~\%) [82.4--100.0] \\ \hline
\texttt{url\_double} & 25 & 25/25 (100.0~\%) [86.7--100.0] & 14/25 (56.0~\%) [37.1--73.3] \\ \hline
\texttt{url\_single} & 25 & 25/25 (100.0~\%) [86.7--100.0] & 25/25 (100.0~\%) [86.7--100.0] \\ \hline
\texttt{whitespace\_tab\_newline} & 20 & 20/20 (100.0~\%) [83.9--100.0] & 20/20 (100.0~\%) [83.9--100.0] \\ \hline
\multicolumn{4}{|l|}{\textit{Transporte}} \\ \hline
\texttt{json} & 68 & 68/68 (100.0~\%) [94.7--100.0] & 63/68 (92.6~\%) [83.9--96.8] \\ \hline
\texttt{multipart} & 68 & 68/68 (100.0~\%) [94.7--100.0] & 62/68 (91.2~\%) [82.1--95.9] \\ \hline
\texttt{urlencoded} & 64 & 64/64 (100.0~\%) [94.3--100.0] & 64/64 (100.0~\%) [94.3--100.0] \\ \hline
\multicolumn{4}{|l|}{\textit{Campo de destino}} \\ \hline
\texttt{bio} & 64 & 64/64 (100.0~\%) [94.3--100.0] & 59/64 (92.2~\%) [83.0--96.6] \\ \hline
\texttt{content} & 68 & 68/68 (100.0~\%) [94.7--100.0] & 63/68 (92.6~\%) [83.9--96.8] \\ \hline
\texttt{title} & 68 & 68/68 (100.0~\%) [94.7--100.0] & 67/68 (98.5~\%) [92.1--99.7] \\ \hline
\multicolumn{4}{|l|}{\textit{Subcadena literal \texttt{alert} en el payload}} \\ \hline
Con la subcadena & 17 & 17/17 (100.0~\%) [81.6--100.0] & 12/17 (70.6~\%) [46.9--86.7] \\ \hline
Sin la subcadena & 183 & 183/183 (100.0~\%) [97.9--100.0] & 177/183 (96.7~\%) [93.0--98.5] \\ \hline
\caption[Cobertura de XSS por estrato, corpus completo]{Cobertura de XSS por estrato sobre los 200 payloads del corpus propio (corrida 1, paquete 0.1.1 en las Configuraciones 2 y 3b)}
\label{tab:anx-xss-estratos}
\end{longtable}
{\footnotesize Nota. Elaboración propia con datos de
\texttt{attack-sim/results\_combined\_config\{2,3a,3b\}\_run1\_20261025.json} y de
\texttt{attack-sim/xss\_corpus\_filtered\_20261025.jsonl} (repositorio
\texttt{logSguarDian-vulnerable-project}); cálculo en
\texttt{docs/evidence/errata-v3-strata-20261025.py}. Las Configuraciones 2 y 3b coinciden en todos
los estratos.}

\begin{longtable}{|L{4.2cm}|L{0.9cm}|L{3.4cm}|L{4.0cm}|}
\hline
\textbf{Estrato} & \textbf{n} & \textbf{Configuraciones 2 y 3b: bloqueados e IC 95~\%} & \textbf{Configuración 3a: bloqueados e IC 95~\%} \\ \hline
\endfirsthead
\multicolumn{4}{l}{\footnotesize\textit{(continuación)}} \\ \hline
\textbf{Estrato} & \textbf{n} & \textbf{Configuraciones 2 y 3b: bloqueados e IC 95~\%} & \textbf{Configuración 3a: bloqueados e IC 95~\%} \\ \hline
\endhead
\multicolumn{4}{r}{\footnotesize\textit{Continúa en la página siguiente}} \\
\endfoot
\endlastfoot
\multicolumn{4}{|l|}{\textit{Contexto de inyección}} \\ \hline
\texttt{attr\_double} & 13 & 13/13 (100.0~\%) [77.2--100.0] & 12/13 (92.3~\%) [66.7--98.6] \\ \hline
\texttt{attr\_single} & 13 & 13/13 (100.0~\%) [77.2--100.0] & 12/13 (92.3~\%) [66.7--98.6] \\ \hline
\texttt{attr\_unquoted} & 9 & 9/9 (100.0~\%) [70.1--100.0] & 8/9 (88.9~\%) [56.5--98.0] \\ \hline
\texttt{html\_body} & 17 & 17/17 (100.0~\%) [81.6--100.0] & 15/17 (88.2~\%) [65.7--96.7] \\ \hline
\texttt{html\_comment} & 3 & 3/3 (100.0~\%) [43.9--100.0] & 3/3 (100.0~\%) [43.9--100.0] \\ \hline
\texttt{js\_string} & 2 & 2/2 (100.0~\%) [34.2--100.0] & 2/2 (100.0~\%) [34.2--100.0] \\ \hline
\texttt{svg} & 4 & 4/4 (100.0~\%) [51.0--100.0] & 4/4 (100.0~\%) [51.0--100.0] \\ \hline
\texttt{url\_href\_src} & 3 & 3/3 (100.0~\%) [43.9--100.0] & 3/3 (100.0~\%) [43.9--100.0] \\ \hline
\multicolumn{4}{|l|}{\textit{Vector de ejecución}} \\ \hline
\texttt{event\_onerror} & 15 & 15/15 (100.0~\%) [79.6--100.0] & 13/15 (86.7~\%) [62.1--96.3] \\ \hline
\texttt{event\_onfocus\_autofocus} & 2 & 2/2 (100.0~\%) [34.2--100.0] & 2/2 (100.0~\%) [34.2--100.0] \\ \hline
\texttt{event\_onload\_body} & 5 & 5/5 (100.0~\%) [56.6--100.0] & 5/5 (100.0~\%) [56.6--100.0] \\ \hline
\texttt{event\_onmouseover} & 14 & 14/14 (100.0~\%) [78.5--100.0] & 12/14 (85.7~\%) [60.1--96.0] \\ \hline
\texttt{event\_ontoggle} & 8 & 8/8 (100.0~\%) [67.6--100.0] & 7/8 (87.5~\%) [52.9--97.8] \\ \hline
\texttt{iframe\_srcdoc} & 1 & 1/1 (100.0~\%) [20.7--100.0] & 1/1 (100.0~\%) [20.7--100.0] \\ \hline
\texttt{javascript\_uri} & 3 & 3/3 (100.0~\%) [43.9--100.0] & 3/3 (100.0~\%) [43.9--100.0] \\ \hline
\texttt{script} & 7 & 7/7 (100.0~\%) [64.6--100.0] & 7/7 (100.0~\%) [64.6--100.0] \\ \hline
\texttt{svg\_onload} & 9 & 9/9 (100.0~\%) [70.1--100.0] & 9/9 (100.0~\%) [70.1--100.0] \\ \hline
\multicolumn{4}{|l|}{\textit{Función ejecutada}} \\ \hline
\texttt{alert} & 12 & 12/12 (100.0~\%) [75.8--100.0] & 10/12 (83.3~\%) [55.2--95.3] \\ \hline
\texttt{confirm} & 8 & 8/8 (100.0~\%) [67.6--100.0] & 8/8 (100.0~\%) [67.6--100.0] \\ \hline
\texttt{document\_cookie} & 5 & 5/5 (100.0~\%) [56.6--100.0] & 5/5 (100.0~\%) [56.6--100.0] \\ \hline
\texttt{eval\_atob} & 8 & 8/8 (100.0~\%) [67.6--100.0] & 8/8 (100.0~\%) [67.6--100.0] \\ \hline
\texttt{fetch\_externo} & 16 & 16/16 (100.0~\%) [80.6--100.0] & 14/16 (87.5~\%) [64.0--96.5] \\ \hline
\texttt{fromcharcode} & 8 & 8/8 (100.0~\%) [67.6--100.0] & 8/8 (100.0~\%) [67.6--100.0] \\ \hline
\texttt{prompt} & 7 & 7/7 (100.0~\%) [64.6--100.0] & 6/7 (85.7~\%) [48.7--97.4] \\ \hline
\multicolumn{4}{|l|}{\textit{Codificación o evasión}} \\ \hline
\texttt{html\_entities\_dec} & 6 & 6/6 (100.0~\%) [61.0--100.0] & 6/6 (100.0~\%) [61.0--100.0] \\ \hline
\texttt{html\_entities\_hex} & 6 & 6/6 (100.0~\%) [61.0--100.0] & 6/6 (100.0~\%) [61.0--100.0] \\ \hline
\texttt{inline\_comment} & 9 & 9/9 (100.0~\%) [70.1--100.0] & 9/9 (100.0~\%) [70.1--100.0] \\ \hline
\texttt{mixed\_case} & 6 & 6/6 (100.0~\%) [61.0--100.0] & 6/6 (100.0~\%) [61.0--100.0] \\ \hline
\texttt{raw} & 5 & 5/5 (100.0~\%) [56.6--100.0] & 5/5 (100.0~\%) [56.6--100.0] \\ \hline
\texttt{slash\_for\_space} & 4 & 4/4 (100.0~\%) [51.0--100.0] & 4/4 (100.0~\%) [51.0--100.0] \\ \hline
\texttt{unicode\_escape} & 7 & 7/7 (100.0~\%) [64.6--100.0] & 7/7 (100.0~\%) [64.6--100.0] \\ \hline
\texttt{url\_double} & 9 & 9/9 (100.0~\%) [70.1--100.0] & 4/9 (44.4~\%) [18.9--73.3] \\ \hline
\texttt{url\_single} & 8 & 8/8 (100.0~\%) [67.6--100.0] & 8/8 (100.0~\%) [67.6--100.0] \\ \hline
\texttt{whitespace\_tab\_newline} & 4 & 4/4 (100.0~\%) [51.0--100.0] & 4/4 (100.0~\%) [51.0--100.0] \\ \hline
\multicolumn{4}{|l|}{\textit{Transporte}} \\ \hline
\texttt{json} & 21 & 21/21 (100.0~\%) [84.5--100.0] & 18/21 (85.7~\%) [65.4--95.0] \\ \hline
\texttt{multipart} & 17 & 17/17 (100.0~\%) [81.6--100.0] & 15/17 (88.2~\%) [65.7--96.7] \\ \hline
\texttt{urlencoded} & 26 & 26/26 (100.0~\%) [87.1--100.0] & 26/26 (100.0~\%) [87.1--100.0] \\ \hline
\multicolumn{4}{|l|}{\textit{Campo de destino}} \\ \hline
\texttt{bio} & 64 & 64/64 (100.0~\%) [94.3--100.0] & 59/64 (92.2~\%) [83.0--96.6] \\ \hline
\caption[Cobertura de XSS por estrato, campo bio]{Cobertura de XSS por estrato sobre los 64 payloads dirigidos al campo \texttt{bio} (corrida 1, paquete 0.1.1 en las Configuraciones 2 y 3b)}
\label{tab:anx-xss-bio}
\end{longtable}
{\footnotesize Nota. Elaboración propia con las mismas fuentes del Cuadro~\ref{tab:anx-xss-estratos}.
Los 5 payloads que la Configuración 3a no bloqueó en \texttt{bio} reciben HTTP 200 de la aplicación.}

\section{Command Injection con separadores codificados y sin codificar}
\label{anx:cmdi-codificacion}

De los 200 payloads de Command Injection, 105 traen al menos un separador ya codificado
(\texttt{\%XX}) y 95 no. El Cuadro~\ref{tab:anx-cmdi-codificacion} los separa por configuración y
por versión del ejecutor. El ejecutor anterior volvía a codificar los separadores ya codificados,
por lo que solo alteraba el primer grupo; el segundo da el mismo resultado con ambos ejecutores.
Con el ejecutor corregido, el WAF sigue bloqueando menos los payloads con separadores codificados
(76.2~\%) que los no codificados (88.4~\%), una diferencia de 12.2 puntos porcentuales (prueba exacta de
Fisher bilateral sobre 80 de 105 frente a 84 de 95, $p = 0.028$; intervalo de confianza de Wald al
95~\% de la diferencia, de 1.8 a 22.6 puntos). Los dos subconjuntos no están controlados por otras
características del payload, por lo que la diferencia no se atribuye a la codificación como causa.
El cuadro conserva los intervalos de Wilson de cada grupo.

\begin{table}[H]
\centering
\small
\begin{tabular}{|L{3.7cm}|L{3.6cm}|L{3.6cm}|L{2.5cm}|}
\hline
\textbf{Configuración y ejecutor} & \textbf{Separadores codificados (n=105): bloqueados e IC 95~\%} & \textbf{Sin codificar (n=95): bloqueados e IC 95~\%} & \textbf{Total (n=200)} \\ \hline
Configuración 1 (sin protección) & 0/105 (0.0~\%) [0.0--3.5] & 0/95 (0.0~\%) [0.0--3.9] & 0/200 (0.0~\%) \\ \hline
Configuración 2 (0.1.1), ejecutor corregido & 95/105 (90.5~\%) [83.4--94.7] & 90/95 (94.7~\%) [88.3--97.7] & 185/200 (92.5~\%) \\ \hline
Configuración 2 (0.1.1), ejecutor anterior (10 de octubre) & 102/105 (97.1~\%) [91.9--99.0] & 90/95 (94.7~\%) [88.3--97.7] & 192/200 (96.0~\%) \\ \hline
Configuración 3a, ejecutor corregido (25 de octubre) & 80/105 (76.2~\%) [67.2--83.3] & 84/95 (88.4~\%) [80.4--93.4] & 164/200 (82.0~\%) \\ \hline
Configuración 3a, ejecutor anterior (10 de octubre; reproducido el 25) & 52/105 (49.5~\%) [40.1--58.9] & 84/95 (88.4~\%) [80.4--93.4] & 136/200 (68.0~\%) \\ \hline
Configuración 3b (0.1.1), ejecutor corregido & 104/105 (99.0~\%) [94.8--99.8] & 93/95 (97.9~\%) [92.6--99.4] & 197/200 (98.5~\%) \\ \hline
\end{tabular}
\caption[Command Injection con separadores codificados y sin codificar]{Command Injection: payloads con separadores ya codificados y sin codificar, por configuración y por versión del ejecutor}
\label{tab:anx-cmdi-codificacion}
\end{table}
{\footnotesize Nota. Elaboración propia con datos de
\texttt{attack-sim/results\_combined\_config\{1,2,3a,3b\}\_run1\_20261025.json},
\texttt{attack-sim/results\_config3a-round5-20261010.json},
\texttt{attack-sim/results\_phase2a\_prefixexecutor\_config3a\_run1\_20261025.json} y los archivos
de resultados de la Configuración~2 del 10 de octubre de 2026; cálculo en
\texttt{docs/evidence/errata-v7b-cmdi-encoding-20261025.py}. Un payload se cuenta como codificado
si contiene al menos un \texttt{\%XX}. La prueba de Fisher y el intervalo de Wald se calcularon en
\texttt{docs/evidence/a3-fisher-codificacion-20261025.py}. Con el ejecutor corregido, 4 payloads de
Command Injection dejaron de bloquearse en la Configuración~3a (\texttt{cmdi\_22},
\texttt{cmdi\_38}, \texttt{cmdi\_52} y \texttt{cmdi\_179}) y 32 pasaron a bloquearse, de modo que el
defecto no sesgaba todos los payloads en la misma dirección.}

\section{Procedencia de la medición unificada del 25 de octubre de 2026}
\label{anx:ronda5-procedencia}

El Cuadro~\ref{tab:anx-ronda5-procedencia} identifica los archivos y las versiones de la medición
de las cuatro configuraciones con el corpus combinado de 800 peticiones. Las rutas son relativas a
la raíz del repositorio \texttt{logSguarDian-vulnerable-project}, rama
\texttt{exp/ronda5-20261025}, sin commits propios.

\begin{table}[H]
\centering
\small
\begin{tabular}{|L{4.2cm}|L{9.3cm}|}
\hline
\textbf{Elemento} & \textbf{Valor} \\ \hline
SHA base de la rama & {\scriptsize\texttt{07e80ead6a64db5db086216322624bc5bfa4edb0}} \\ \hline
Candidatos de XSS (400), SHA-256 & {\scriptsize\texttt{98d6e846899feedb0e736f3cb842d13c\allowbreak{}8d506751fb60911318bf840137883614}} \\ \hline
Corpus de XSS (200), SHA-256 & {\scriptsize\texttt{ef67ffd34d09a59c4cf8212722f78d5a\allowbreak{}2e8894f8d6a38d6e63a2f4598b6fae0d}} \\ \hline
Corpus combinado (800), SHA-256 & {\scriptsize\texttt{f1fb0b0a2c552636358785ec323aa10c\allowbreak{}76519638d0da1cb62ce131be2330263d}} \\ \hline
Protocolo de XSS, SHA-256 & {\scriptsize\texttt{ae52abc3227feb00650e98e0e1e81185\allowbreak{}fc76a56d8bd6b93ba1ca5cb3cc53134f}} \\ \hline
Manifiesto de congelación de XSS (versión final), SHA-256 & {\scriptsize\texttt{b3ccb69cc6329b520d4b0e91e28e9c10\allowbreak{}30940afedbdfe3915408ca31faee0496}} \\ \hline
Ejecutor corregido (\texttt{run\_large\_corpus.py}, commit \texttt{00ea3a3}), SHA-256 & {\scriptsize\texttt{422393460dfea808efe411d358166d37\allowbreak{}b662bbcaad77843ede0247c684c61af1}} \\ \hline
Ejecutor anterior (\texttt{00ea3a3\^{}}), SHA-256 & {\scriptsize\texttt{0c84de551f1138bbeb1ef7a9ed1801d8\allowbreak{}3e4c52a7312c04b873a75de57a286623}} \\ \hline
Emisor de XSS (\texttt{send\_xss\_corpus.py}), SHA-256 & {\scriptsize\texttt{18c645527d1c6944734edb899ce10279\allowbreak{}928285abef1a627691c06223dcaa75ee}} \\ \hline
Paquete \texttt{logsguardian@0.1.1}, SHA-256 del archivo & {\scriptsize\texttt{abf8d5c47036d9c3a289bc36158f1caa\allowbreak{}a91b9504689aa97193e51c6808258d15}} \\ \hline
Modelo \texttt{rf.onnx} (par~C), SHA-256 & {\scriptsize\texttt{25b407e649283606a231665a5a8c8ea1\allowbreak{}c2c7ad4fd0374a8945b3f119651b0870}} \\ \hline
Modelo \texttt{if.onnx} (par~C), SHA-256 & {\scriptsize\texttt{5be28ee8437643b2f13708eee1ac3175\allowbreak{}348a350ea6fb8f8daeb09bd5717f797e}} \\ \hline
Imagen del WAF, digest & {\scriptsize\texttt{sha256:\allowbreak f12da37e72e29d1f06e7e829780f1016\allowbreak{}606d0799eac2966e58ef3819a3d30a50}} \\ \hline
Versiones del WAF & Nginx 1.30.5, ModSecurity 3.0.17, OWASP CRS 4.30.0; paranoia 1; \texttt{SecRuleEngine On}; umbral de anomalía de entrada 5; sin exclusiones activas \\ \hline
\texttt{Dockerfile} de la aplicación, SHA-256 & {\scriptsize\texttt{7f1edcdd31b1bd63b3d84d71a49f91b6\allowbreak{}a2cb63c8435413b6bab3e4e09d84203d}} \\ \hline
Resultados de la Configuración~1, corrida 1, SHA-256 & {\scriptsize\texttt{d88f968d203762aa469998968a798010\allowbreak{}18992f444c8ad7ff1bfef65986610903}} \\ \hline
Resultados de la Configuración~2, corrida 1, SHA-256 & {\scriptsize\texttt{577e6e002fb4989c719562a63a4db064\allowbreak{}82e40bff46ad68eb1faefcbb3b9b59aa}} \\ \hline
Resultados de la Configuración~3a, corrida 1, SHA-256 & {\scriptsize\texttt{1cf8badbbcc7f92f633aabd07d5d2083\allowbreak{}55d94a97802587d510ff7f51818b502c}} \\ \hline
Resultados de la Configuración~3b, corrida 1, SHA-256 & {\scriptsize\texttt{1cfd0b5ec657d7fdc6816d10eb2b3f89\allowbreak{}92cf687d3c22d0a0d957db3b1048ba8e}} \\ \hline
\end{tabular}
\caption[Procedencia de la medición unificada]{Procedencia de la medición unificada de las cuatro configuraciones del 25 de octubre de 2026}
\label{tab:anx-ronda5-procedencia}
\end{table}
{\footnotesize Nota. Elaboración propia con datos de \texttt{docs/thesis-delta-xss.md},
\texttt{docs/thesis-delta-configs.md} y \texttt{docs/thesis-delta-errata-20261025.md}; las sumas
se recalcularon el 25 de octubre de 2026. El manifiesto de congelación es el de su versión
final: incluye los recálculos posteriores a la medición, y la garantía de congelación descansa en
los hashes de los archivos congelados que contiene, no en el del manifiesto. La etiqueta de la
imagen del WAF es flotante y hoy resuelve a otro manifiesto, por lo que se cita el digest medido. La imagen local de la aplicación
tuvo el identificador \texttt{a75c14208ccd}; su digest no se verificó en cada arranque.}

\section{Sensibilidad de la Configuración 2 a las cabeceras del proxy}
\label{anx:sensibilidad}

El Cuadro~\ref{tab:anx-sensibilidad} compara la Configuración~2 con el paquete 0.1.1 en dos
condiciones, ambas sobre el corpus combinado de 800 peticiones y con el ejecutor corregido: sin
cabeceras extra (corrida 1 de la medición unificada) y con las cuatro cabeceras
\texttt{X-Real-IP}, \texttt{X-Forwarded-For}, \texttt{X-Forwarded-Port} y
\texttt{X-Forwarded-Proto} que agrega el proxy, con los valores de la captura de la Fase 2b
(192.168.65.1, 192.168.65.1, 8080 y http). Es una sensibilidad: las cifras principales del
Capítulo~\ref{cap:resultados} no cambian. La condición sin cabeceras se repitió en el mismo
contenedor y reprodujo la corrida 1 sin diferencias en ningún payload.

\begin{table}[H]
\centering
\small
\begin{tabular}{|L{2.7cm}|L{3.7cm}|L{3.7cm}|L{1.6cm}|L{1.6cm}|}
\hline
\textbf{Categoría} & \textbf{Sin cabeceras: bloqueados e IC 95~\%} & \textbf{Con cuatro cabeceras: bloqueados e IC 95~\%} & \textbf{A bloqueado} & \textbf{A no bloqueado} \\ \hline
SQLi & 199/200 (99.5~\%) [97.2--99.9] & 200/200 (100.0~\%) [98.1--100.0] & 1 & 0 \\ \hline
Path Traversal & 200/200 (100.0~\%) [98.1--100.0] & 200/200 (100.0~\%) [98.1--100.0] & 0 & 0 \\ \hline
Command Injection & 185/200 (92.5~\%) [88.0--95.4] & 194/200 (97.0~\%) [93.6--98.6] & 11 & 2 \\ \hline
XSS & 200/200 (100.0~\%) [98.1--100.0] & 200/200 (100.0~\%) [98.1--100.0] & 0 & 0 \\ \hline
\textbf{Total} & 784/800 (98.0~\%) [96.8--98.8] & 794/800 (99.2~\%) [98.4--99.7] & 12 & 2 \\ \hline
\end{tabular}
\caption[Sensibilidad de la Configuración 2 a las cabeceras del proxy]{Sensibilidad de la Configuración 2 (paquete 0.1.1) a las cuatro cabeceras que agrega el proxy, sobre 800 peticiones}
\label{tab:anx-sensibilidad}
\end{table}
{\footnotesize Nota. Elaboración propia con datos de
\texttt{attack-sim/results\_combined\_config2\_run1\_20261025.json} (sin cabeceras) y de
\texttt{attack-sim/results\_s1\_config2\_con\_cabeceras\_run1\_20261025.json} (con cabeceras); cálculo
en \texttt{docs/evidence/s1-analisis-20261025.py}. Una corrida por condición. Pasan a bloqueado: \texttt{cmdi\_112}, \texttt{cmdi\_120}, \texttt{cmdi\_134}, \texttt{cmdi\_155}, \texttt{cmdi\_28}, \texttt{cmdi\_31}, \texttt{cmdi\_36}, \texttt{cmdi\_38}, \texttt{cmdi\_40}, \texttt{cmdi\_66}, \texttt{cmdi\_7}, \texttt{sqli\_75}. Dejan de bloquearse: \texttt{cmdi\_51}, \texttt{cmdi\_94}, órdenes con \texttt{sleep} que llegan a la aplicación y terminan por tiempo de
espera, lo que se cuenta como no detectado. Prueba exacta de McNemar sobre los 14 payloads
discordantes (12 frente a 2): $p = 0.013$. De los 13 payloads que la Configuración~3b bloquea y la
Configuración~2 sin cabeceras no, la Configuración~2 con cabeceras bloquea 12 (todos salvo
\texttt{cmdi\_9}); la unión de la Configuración~2 con cabeceras y la 3a es de 797 payloads, igual que
la 3b, y el WAF suma 3 (\texttt{cmdi\_9}, \texttt{cmdi\_51} y \texttt{cmdi\_94}) a la librería con
cabeceras. La imagen local de la aplicación de esta sensibilidad (\texttt{sha256:d9b12672f000...}) se
reconstruyó porque cambió el contexto de build; el paquete 0.1.1 y los modelos se verificaron por
hash dentro del contenedor.}

\ifdefined\modoensamble
\else
	\section*{Referencias}
	\printbibliography[heading=none]
	\end{document}
\fi
